\section{Por qué la Embodied Intelligence es más difícil que conversar}

En el mundo virtual, un error suele manifestarse como un texto equivocado; en el mundo físico, un error puede alterar objetos, dañar el hardware e incluso lastimar a personas. Por eso la inteligencia corporizada no solo necesita reconocer «qué es esto», sino también predecir «qué ocurrirá después de que yo haga esta acción». Esto amplía el problema de la comprensión semántica estática a la toma de decisiones en lazo cerrado y al modelado de consecuencias.

\subsection{Los casos de falla revelan una carencia de world model}

Esta sección parte de dos casos de falla para determinar qué capacidades deberán aportar los modelos posteriores. El ponente muestra un ejemplo en el que un robot, al colocar una lata de refresco de cola en el fregadero, abre por accidente la llave del agua, y otro en el que un brazo robótico, al retraerse, voltea un recipiente. Estos errores no son un simple object detection failure; es posible que el robot haya visto la lata, el fregadero y el brazo robótico, pero no modeló lo suficiente el efecto de la secuencia de acciones sobre el líquido, la gravedad, las colisiones y su propia postura.

\lecturefigure{slide-02-why-embodied-intelligence.jpg}{La complejidad del mundo real y el objetivo de la ambient intelligence}{Diapositiva V002 recuperada de la grabación oficial; 00:04:06}

Esta carencia puede expresarse con una ecuación mínima de transición de estados:
$$
s_{t+1}=f(s_t,a_t,\xi_t),\qquad o_t=h(s_t)+\epsilon_t.
$$
Aquí $s_t$ es el estado real del entorno en el instante $t$, $a_t$ es la acción del robot, $f$ es la transición de estados afectada por la dinámica y el contacto, y $\xi_t$ representa las perturbaciones no modeladas; $o_t$ es la observation del sensor, $h$ es la función de observación y $\epsilon_t$ es el ruido de medición. El robot solo ve $o_t$, pero debe inferir el estado oculto $s_t$ y predecir las consecuencias de sus acciones.

\begin{knowledgebox}{First use: world model}
\term{World model (modelo del mundo)} es una representación interna de la relación entre el estado del entorno, las acciones y los resultados futuros. Puede ser un modelo dinámico explícito, una latent transition implícita o un predictor de video, y también puede estar codificado en parte dentro de la policy. Su valor no consiste en replicar el mundo entero, sino en sustentar la predicción, la planeación y la detección de errores relevantes para la tarea.
\end{knowledgebox}

\teachervoice{Fei Xia usa la conducta inesperada de «el robot coloca la lata en el fregadero, pero abre la llave del agua» para recordar a los estudiantes que el mundo real produce consecuencias de las acciones que quienes entrenan el modelo no previeron. Este ejemplo no es un blooper simpático: muestra que, si la policy carece de un consequence model, incluso un error aleatorio puede llevar a un estado peligroso.}

\subsection{Gibson: la inteligencia surge del lazo cerrado entre el agent y el entorno}

Las fallas anteriores conducen de forma natural a la ecological view de James J. Gibson. En lugar de preguntar solo qué almacena el modelo internamente, conviene preguntar en qué entorno se coloca al agent, qué acciones puede ejecutar y qué retroalimentación recibirá. Los niños aprenden la gravedad, la fricción, el soporte, el alcance y los usos de los objetos mediante años de interacción corporal; si un robot no tiene una experiencia equivalente, difícilmente adquirirá el mismo sentido común con solo unas cuantas teleoperation trajectory.

\lecturefigure{slide-03-gibson-ecological-perception.jpg}{La percepción ecológica de Gibson: comprender el entorno en el que se encuentra la head}{Diapositiva V003 recuperada de la grabación oficial; 00:05:18}

\begin{knowledgebox}{First use: affordance}
\term{Affordance (posibilidad de acción)} describe lo que el entorno permite hacer a un agent determinado; por ejemplo, una manija «se puede sujetar», una superficie plana «permite colocar cosas», un botón «se puede presionar». No es una propiedad puramente visual del objeto, sino que depende de la morfología del robot, de su espacio de acciones y de la tarea. Un mismo vaso ofrece affordance distintas a una pinza de dos dedos, a una mano humana y a una base móvil.
\end{knowledgebox}

\subsection{La simulation amplifica la experiencia, pero no sustituye al mundo real}

Una de las implementaciones de ingeniería del playpen seguro es la simulation. El ponente repasa cómo, en trabajos como iGibson, construyó escenas interiores a gran escala, renderizó la observation y dejó que el agent aprendiera en entornos interactivos. Las ventajas de la simulation son que es repetible, paralelizable, que permite acceder al estado oculto y generar label automáticamente; su riesgo es que existe una gap con el mundo real en lo visual, en la dinámica, en el contacto y en la distribución de tareas.

\lecturefigure{slide-04-simulation-digital-twin-perception.jpg}{La simulation proporciona a la vez rendered observation, physics y semantic labels}{Diapositiva V004 recuperada de la grabación oficial; 00:06:36}

\lecturefigure{slide-05-large-realistic-scenes.jpg}{Realistic scenes a gran escala: de los recursos geométricos a los entornos interactivos}{Diapositiva V005 recuperada de la grabación oficial; 00:07:24}

\begin{warningbox}{La simulation no equivale automáticamente a un world model}
El simulador contiene la dinámica y la distribución de recursos que escribieron sus desarrolladores; que un agent tenga éxito en el simulador puede deberse solo a que aprovechó un shortcut en las texturas, las colisiones o el mecanismo de reset. Al leer un simulation result hay que distinguir si el entorno cubre la variación real, si el ruido de los sensores es razonable, si el modelo de contacto es confiable y si la policy se validó de forma independiente en un robot real.
\end{warningbox}

\subsection{Las tareas de horizonte largo exponen los cuellos de botella de datos y de composicionalidad}

La sujeción en un solo paso puede aprenderse con una gran cantidad de datos repetidos, pero las tareas domésticas de horizonte largo exigen mantener el estado a lo largo de múltiples room, object, failure recovery y subgoal. La household demo de la figura no pretende demostrar que el robot doméstico ya está resuelto, sino mostrar que el pipeline tradicional de imitation learning se encarece rápidamente en escala de escenas, número de interacciones y combinación de tareas.

\lecturefigure{slide-06-long-horizon-home-tasks.jpg}{Tareas domésticas de horizonte largo: cadenas de acciones, cambios de estado y recuperación ante fallas}{Diapositiva V006 recuperada de la grabación oficial; 00:08:15}

\lecturefigure{slide-07-internet-ai-vs-embodied-ai.jpg}{Internet AI y Embodied AI: comparten módulos semánticos, pero difieren en el lazo cerrado de acción}{Diapositiva V007 recuperada de la grabación oficial; 00:09:12}

\lecturefigure{slide-08-simulation-to-foundation-models.jpg}{De 50k episodes / 1.25M interactions hacia los Foundation Models}{Diapositiva V008 recuperada de la grabación oficial; 00:11:15}

\teachervoice{El ponente describe su paso de la large-scale simulation a los foundation models como un cambio de línea de investigación: solo ampliar escenas e interacciones sigue siendo insuficiente para cubrir el mundo real, mientras que los modelos web-scale pueden aportar priors sobre objetos, lenguaje y sentido común. La palabra central aquí es «complementar el prior», no «la simulation y los robot data ya no son necesarios».}

\subsection{Resumen de la sección}

Las dificultades de la inteligencia corporizada provienen del estado parcialmente observable, las acciones continuas, las consecuencias físicas y la composición de horizonte largo. La simulation puede amplificar la experiencia, pero sigue existiendo la sim-to-real gap; la oportunidad de los foundation models está en incorporar conocimiento previo visual y lingüístico a gran escala, y la siguiente sección debe responder cómo se conecta ese conocimiento previo con el control de bajo nivel.
